\documentclass[CJK]{beamer}
\usepackage{CJKutf8}
\usepackage{beamerthemesplit}
\usetheme{Malmoe}
\useoutertheme[footline=authortitle]{miniframes}
\usepackage{amsmath}
\usepackage{amssymb}
\usepackage{graphicx}
\usepackage{color}
\usepackage{slashed}
\usepackage{simplewick}
\graphicspath{{../figures/}}
\def\be{\begin{equation}}
\def\ee{\nonumber\end{equation}}
\def\bea{\begin{eqnarray}}
\def\eea{\nonumber\end{eqnarray}}
\def\ii{{\dot{\imath}}}
\def\bch{\begin{CJK}{UTF8}{gbsn}}
\def\ech{\end{CJK}}
\def\bex{\begin{minipage}{0.3\textwidth}\includegraphics[width=1in]{jugelizi.png}\end{minipage}\begin{minipage}{0.6\textwidth}}
\def\eex{\end{minipage}}
\def\chtitle#1{\frametitle{\bch#1\ech}}
\def\skipline{{\vskip0.1in}}
\def\skiplines{{\vskip0.2in}}
\def\lagr{{\mathcal{L}}}
\def\hamil{{\mathcal{H}}}
\def\vecv{{\mathbf{v}}}
\def\vecx{{\mathbf{x}}}
\def\veck{{\mathbf{k}}}
\def\vecp{{\mathbf{p}}}
\def\vecn{{\mathbf{n}}}
\def\vecA{{\mathbf{A}}}
\def\vecP{{\mathbf{P}}}
\def\vecsigma{{\mathbf{\sigma}}}
\def\hatJn{{\hat{J_\vecn}}}
\def\hatJx{{\hat{J_x}}}
\def\hatJy{{\hat{J_y}}}
\def\hatJz{{\hat{J_z}}}
\def\hatj#1{\hat{J_{#1}}}
\def\hatphi{{\hat{\phi}}}
\def\hatq{{\hat{q}}}
\def\hatpi{{\hat{\pi}}}
\def\vel{\upsilon}
\def\Dint{{\mathcal{D}}}
\def\adag{{\hat{a}^\dagger}}
\def\bdag{{\hat{b}^\dagger}}
\def\cdag{{\hat{c}^\dagger}}
\def\ddag{{\hat{d}^\dagger}}
\def\hata{{\hat{a}}}
\def\hatb{{\hat{b}}}
\def\hatc{{\hat{c}}}
\def\hatd{{\hat{d}}}
\def\hatN{{\hat{N}}}
\def\hatH{{\hat{H}}}
\def\hatp{{\hat{p}}}
\def\Fup{{F^{\mu\nu}}}
\def\Fdown{{F_{\mu\nu}}}
\def\newl{\nonumber \\}
\def\SIkm{\mathrm{km}}
\def\SIyr{\mathrm{yr}}
\def\SIGyr{\mathrm{Gyr}}
\def\SIeV{\mathrm{eV}}
\def\SIGeV{\mathrm{GeV}}
\def\SIm{\mathrm{m}}
\def\SIcm{\mathrm{cm}}
\def\SIJ{\mathrm{J}}
\def\SIs{\mathrm{s}}
\def\SIkg{\mathrm{kg}}
\def\SIg{\mathrm{g}}
\def\vece{\mathrm{e}}
\def\bmat#1{\left(\begin{array}{#1}}
\def\emat{\end{array}\right)}
\def\bcase#1{\left\{\begin{array}{#1}}
\def\ecase{\end{array}\right.}
\def\calM{{\mathcal{M}}}
\def\calT{{\mathcal{T}}}
\def\calR{{\mathcal{R}}}
\def\barpsi{\bar{\psi}}
\def\baru{\bar{u}}
\def\barv{\bar{\upsilon}}
\def\bmini#1{\begin{minipage}{#1\textwidth}}
\def\emini{\end{minipage}}
\def\qeq{\stackrel{?}{=}}
\def\torder#1{\mathcal{T}\left(#1\right)}
\def\rorder#1{\mathcal{R}\left(#1\right)}


\title{Quantum Field Theory I \\ Homework 3}
  \author{}
  \date{}


\begin{document}

\begin{frame}
 
\begin{center}
\begin{Large}
\bch
量子场论 I 

{\vskip 0.3in}

第三次课后作业 （共八次，每次2.5分）

交作业时间: 10月24日，星期一，13：30pm

\ech
\end{Large}
\end{center}

\vskip 0.2in

\bch
课件下载
\ech
https://github.com/zqhuang/SYSU\_QFTI

\end{frame}

\begin{frame}
\chtitle{第1题(0.5分）}
\bch
证明自由$U(1)$规范场$A^\mu$在库仑规范$A^0=0, \nabla\cdot \vecA=0$下的Hamilton密度为
$$\hamil = \frac{1}{2}|\dot\vecA|^2 + \frac{1}{2}|\nabla\times \vecA|^2$$
并证明在傅立叶空间Hamilton量可以写成
$$H = \int d^3\veck\left[\frac{1}{2}|\dot\vecA|^2 + \frac{1}{2}|\veck \times \vecA|^2\right]$$
\ech
\end{frame}

\begin{frame}
\chtitle{第2题(0.5分）}
\bch
利用上题的Hamilton量的表达式以及我们在课上得到的$\hat{\vecA}$在傅立叶空间的解：
$$\hat{\vecA} = \frac{1}{\sqrt{2|\veck|d^3\veck}}\sum_{s=\pm 1} \vece_{\veck,s}\left(\hata_{\veck, s} + \adag_{-\veck, s}\right)$$
证明Hamilton量的算符表达式为
$$\hat{H} = \sum_{\veck} \sum_{s=\pm 1}(\hat{N}_{\veck,s}+1/2)|\veck|$$
其中 $\hat{N}_{\veck, s} \equiv \adag_{\veck,s}\hata_{\veck, s}$是动量为$\veck$，自旋为$s$的粒子数算符。
\ech
\end{frame}

\begin{frame}
\chtitle{第3题(0.5分）}
\bch
证明自由$U(1)$规范场$A^\mu$在库仑规范$A^0=0, \nabla\cdot \vecA=0$下的动量
$$P^i = \int d^3\vecx\, (\partial^0\vecA \cdot \partial^i\vecA)$$
是守恒量。并把它写成傅立叶空间的积分。
\skipline
\skipline

提示：用Nother定理和三个方向的空间平移不变性。 
\ech
\end{frame}


\begin{frame}
\chtitle{第4题(0.5分）}
\bch
利用上题的动量$P^i$的表达式以及我们在课上得到的$\hat{\vecA}$在傅立叶空间的解：
$$\hat{\vecA} = \frac{1}{\sqrt{2|\veck|d^3\veck}}\sum_{s=\pm 1} \vece_{\veck,s}\left(\hata_{\veck, s} + \adag_{-\veck, s}\right)$$
证明动量$\vecP \equiv (P^1, P^2, P^3)$的算符表达式为
$$\hat{\vecP} = \sum_{\veck}\sum_{s = \pm 1} \veck \hat{N}_{\veck, s}$$
其中 $\hat{N}_{\veck, s} \equiv \adag_{\veck,s}\hata_{\veck, s}$是动量为$\veck$，自旋为$s$的粒子数算符。
\ech
\end{frame}


\begin{frame}
\chtitle{第5题(0.5分）}
\bch
对旋量$\psi$，证明$\bar\psi\slashed{\partial}\psi$是洛仑兹变换下的标量。
\ech
\end{frame}


\end{document}
